%!TEX root = ./../../springer-mono-format/book.tex
%!TEX root = ./../../thesis-trier.tex

Here, in \cref{prelim:sec:graphs:notions}, we introduce undirected and directed graphs, 
and related notions. \cref{prelim:sec:graphs:conn} 
investigates different connectedness notions, paths, and cycles.
In \cref{prelim:sec:graphs:funcgraph} we
introduce the graph of a function and of a homogeneous relation.

\subsection{Basic Notions}
\label{prelim:sec:graphs:notions}

\IN{edge of a graph}
\IN{vertex of a graph} 
\index{node of a graph|see{vertex of a graph}}
\IN{directed graph}
\index{digraph|see{directed graph}}
\IN{endpoints of edge}
\IN{outgoing edge} 
\IN{incoming edge}
\IN{starting vertex}
\IN{initial vertex} 
\IN{source vertex} 
\IN{ending vertex}
\IN{final vertex}
\IN{terminal vertex}
\IN{outdegree of vertex}
\IN{indegree of vertex}
An \emph{undirected graph}, or just \emph{graph}, (resp.~a \emph{directed graph}, or \emph{digraph})
is a tuple $G = (V, E)$ where~$V$ is the set of \emph{vertices} (or \emph{nodes}) 
and $E \subseteq \{ \{u,v\} \mid u,v \in V \}$
(resp.~$E \subseteq V \times V$)
the set of \emph{edges}.\footnote{Here, $u = v$ (a self-loop) is permitted, which some authors exclude from the definition.} 
For an edge $\{ u,v \} \in E$ (resp.~$(u,v) \in E$)
we call~$u$ and~$v$ its \emph{endpoints}. 
In case of digraphs, $(u,v)$ is called an \emph{outgoing edge w.r.t.~$u$} (and~$u$ is called the \emph{starting/initial/source vertex})
and an \emph{incoming edge w.r.t.~$v$} (and~$v$ the \emph{ending/final/terminal vertex}).
For finite~$V$ and $v \in V$ 
the \emph{outdegree} (resp.~\emph{indegree})
of $v$ is $|\{ u \in V : (v,u) \in E \}|$ (resp.~$|\{ u \in V : (u,v) \in E \}|$).

\IN{subgraph}
\IN{automorphism!of graphs}
\IN{automorphism group!of graphs}
\label{def:AutG}
A \emph{subgraph} of an undirected (resp.~directed) graph $G = (V,E)$
is an undirected (resp.~directed) graph $G = (V', E')$
such that $V' \subseteq V$ and $E' \subseteq E$.
A \emph{(graph) morphism} of two undirected (resp.~directed)
graphs $G = (V,E)$ and $G' = (V', E')$
is a map $\varphi \colon V \to V'$
such that, for all $u,v \in V$,
\[
\{ u,v \} \in E \Rightarrow \{ \applyfunc{u}{\varphi}, \applyfunc{v}{\varphi} \} \in E'
\quad (\text{resp. } (u,v) \in E \Rightarrow (\applyfunc{u}{\varphi}, \applyfunc{v}{\varphi}) \in E').
\] 
A \emph{(graph) automorphism} (for undirected and directed graphs)
is a morphism $\varphi \colon V \to V$ 
of a graph onto itself such that there exists a morphism $\psi \colon V \to V$
with $\concat{\varphi}{\psi} = \concat{\psi}{\varphi} = \operatorname{id}_V$.
These are precisely those morphisms $\varphi \colon V \to V$ such
that 
\[
\{ u,v \} \in E \Leftrightarrow \{ \applyfunc{u}{\varphi}, \applyfunc{v}{\varphi} \} \in E'
\quad (\text{resp. } (u,v) \in E \Leftrightarrow (\applyfunc{u}{\varphi}, \applyfunc{v}{\varphi}) \in E').
\] 
The set of all automorphisms forms a permutation group, called 
the \emph{automorphism group} $\operatorname{Aut}(G)$ of the graph.

Automorphisms of finite graphs are precisely the bijective morphisms from~$V$ onto~$V$
since in this case $\varphi(E) \subseteq E$ and $|\varphi(E)| = |E|$ 
imply $\varphi(E) = E$. However, this fails for infinite graphs,
as shown by the following example.

\begin{example}
	$G = (\mathbb Z, E)$ with $E = \{ \{u,v\} \mid \{n,n+1\} \text{ and } n > 0 \}$.
	The map $\varphi \colon \mathbb Z \to \mathbb Z$
	with $\varphi(n) = n + 1$ for $n \in \mathbb Z$
	is a morphism, but it is not an automorphism since
	$\{0,1\} \notin E$ and $\varphi(\{0,1\}) = \{1,2\} \in E$.
\end{example}

\IN{symmetric digraph}
\IN{oriented digraph}
Let $G = (V,E)$ be a digraph. We call $G$ \emph{symmetric} (resp.~\emph{oriented})
if $(u,v) \in E$ implies $(v,u) \in E$ (resp.~$(v,u) \notin E$) for $u \ne v$.

We often give a graph (resp.~digraph) in graphical form by drawing the vertices
as circles and the edges as connecting lines (resp.~connecting arrows)
between them.

Groups are also often denoted by $G$, but the context should make
it clear whether we refer to a graph or a group.

\subsection{Connectedness Notions} 
\label{prelim:sec:graphs:conn} 

Here, we introduce various notions related to connectedness in graphs. Let $k \ge 1$. 

\begin{enumerate}	
	\IN{walk in a graph}
	\IN{cycle in a graph}
	\IN{path in a graph}
	\item A \emph{walk} (of length~$n \ge 1$)\footnote{There exist no walks of length zero according to our definition.}
	in an undirected (resp.~directed) graph $G = (V,E)$ from $u \in V$ (\emph{start}) to $v \in V$ (\emph{end}) is a sequence
	of edges $e_1, \ldots, e_n$
	with $e_i = \{ v_{i-1}, v_i\}$ (resp.~$e_i = (v_{i-1}, v_i)$) for all $i \in \{1,\ldots,n\}$
	such that $v_0 = u$ and $v_n = v$.
	We call the walk a \emph{path}		
	(resp.~\emph{cycle} (at $u$))
	if $v_0, v_1, \ldots, v_n$ are distinct (resp.~$v_0, v_1,\ldots, v_{n-1}$ are distinct and $v_0 = v_n$).\footnote{The 
		terminology varies in the literature: some authors use path where we use walk, and reserve simple path for a walk whose vertices are all distinct.}
	Often, when the specific edges involved do not matter, we refer to a walk/path/cycle by its sequence of vertices
	$v_0, v_1, \ldots, v_n \in V$.\footnote{Paths and cycles can also be identified
		with graphs and digraphs such that a graph or digraph
		contains a path if and only if it contains these as subgraphs.
		However, this does not work for walks since they can have repeated vertices
		or edges.}		
	
\IN{acyclic graph}
\IN{directed acyclic graph}
\IN{forest}
\index{DAG|see{directed acyclic graph}}
\item A directed or undirected graph $G$ is called \emph{acyclic}
if there exists no cycle in~$G$. Equivalently, every walk is a path.
For short, a directed acyclic graph is called a \emph{DAG}. An undirected acyclic graph
is also called a \emph{forest}.

\IN{$k$-connected graph}
\IN{strongly $k$-connected graph}
\IN{connected graph} \IN{strongly connected!digraph}
\item An undirected (resp.~directed) graph $G = (V, E)$ is \emph{$k$-connected} (resp.~\emph{strongly $k$-connected})
if for distinct vertices $u,v \in V$ and every $(k-1)$-set $A \subseteq V$ there exists
a path $u_0, u_1, \ldots, u_n$ from $u$ to $v$
such that $\{u_1, \ldots, u_{n-1}\} \cap A = \emptyset$.
For $k = 1$ we also call $G$ just \emph{connected} (resp.~\emph{strongly connected}).
	
	\IN{strongly $\frac{3}{2}$-connected digraph}
	\item We call the digraph $G = (V,E)$ \emph{strongly $\frac{3}{2}$-connected}
	if it is strongly connected and for every path $v_0, v_1, \ldots, v_n$
	and $0 < i < n$ there exists a path $u_0, u_1, \ldots, u_m$
	with $\{u_0, u_m \} = \{ v_0, v_n \}$
	and $v_i \notin \{u_1, \ldots, u_{m-1}\}$. Equivalently, $G$ is strongly connected and
	for every three distinct $u,v,w \in V$
	there exists a path from~$u$ to~$v$ or from $v$ to $u$ not traversing~$w$, i.e.,
	we find a path avoiding $w$.
	
	\IN{weakly $k$-connected digraph}
	\item A digraph $G = (V, E)$ is called \emph{weakly $k$-connected}
	if the undirected graph with vertex set $V$ and edge set $\{ \{u,v \} \mid (u,v) \in E \}$
	is $k$-connected. A weakly $1$-connected graph is simply called \emph{weakly connected}.
	
	\IN{unilaterally connected!digraph} \index{semiconnected graph|see{unilatterally connected graph}}
	\item A digraph $G = (V, E)$ is called \emph{unilaterally connected} (or \emph{semiconnected})
	if for distinct vertices $u, v \in V$ there exists a path from $u$ to $v$ or from $v$ to $u$.
		
	\IN{connected component}
	\IN{strongly connected component!in digraph}
	\IN{weakly connected component}  
	\item In an undirected (resp.~directed) graph $G = (V, E)$
    a \emph{connected component} (resp.~\emph{strongly connected component}) 
    is a subset $S \subseteq V$ maximal for subset inclusion
    such that there exists a path from~$u$ to~$v$
    for distinct $u,v \in S$.
	A subset of a digraph $G = (V,E)$ 
	is a \emph{weakly connected component} if 
	it is a connected component of the undirected graph with vertex set $V$
	and edge set $\{ \{u,v\} \mid (u,v) \in E \}$.
\end{enumerate} 

\IN{connectivity}
\begin{remark} \label{prelim:rem:connectedconnectivity}
	The \emph{connectivity}
	of a finite graph (resp.~digraph) with $n$ vertices
	in which not all vertices are joined by an edge
	is the minimal number of vertices whose
	deletion results in a
	graph (resp.~digraph)
	that is not connected (resp.~not strongly connected).
    This also equals the largest~$k$
    such that the graph (resp.~digraph)
    is $k$-connected (resp.~strongly $k$-connected).
        
    If all vertices are joined by an edge,
    deletion of any number of vertices never
    destroys connectedness or strong connectedness
    and it is $k$-connected or strongly $k$-connected
    for all $k \ge 1$.
    In this case, the connectivity is defined to be $n - 1$.
\end{remark}

\IN{topological sorting!of a DAG}
In a DAG $G = (V, E)$,
when putting $u,v \in V$
in relation if there exists a path from $u$ to $v$, we obtain 
a partial order. 
When saying we \emph{topologically sort the DAG}
we refer to a topological sorting of the vertices
w.r.t.~this partial order.

The following classical result 
was proved by Menger for undirected graphs~\cite{Menger1927}
and extended to digraphs by Dirac~\cite{dirac1953extensions}.
A short proof for digraphs
is given in~\cite{DBLP:journals/dm/Goring00}.

\begin{theorem}[Menger-Dirac~\cite{dirac1953extensions,Menger1927}]
An undirected (resp.~directed) graph $G = (V,E)$ is \emph{$k$-connected} (resp.~\emph{strongly $k$-connected})
if for distinct $u,v \in V$ there exist at least~$k$
paths from $u$ to $v$ such that for two such paths $v_0, v_1, \ldots, v_n$
and $u_0, u_1, \ldots, u_m$ with  $v_0 = u_0 = u$, $v_n = u_m = v$
we have $\{ v_1, \ldots, v_{n-1}\} \cap \{ u_1, \ldots, u_{m-1}\} = \emptyset$.
\end{theorem}

In the context of the present work, only strong $\frac{3}{2}$-connectivity
and (strong) $k$-connectivity with $k \le 2$ will be of interest.
The implications from \cref{prelim:fig:connectivity-notions} 
are easily seen. Figure~\ref{prelim:fig:connectedness-notions}
shows that the converse implications do not hold. 
Furthermore, it also gives an example of a strongly connected digraph
that is weakly $2$-connected but not strongly~$\frac{3}{2}$-connected; thus
strong connectivity and weak $2$-connectivity do not imply 
strong~$\frac{3}{2}$-connectivity.

\begin{figure}[htb]
\begin{center} 
	\begin{tikzcd}[arrows=Rightarrow]
		&                                          & \textrm{weakly $2$-connected} \\
		\textrm{strongly $2$-connected} \arrow[r]  & \textrm{strongly $\frac{3}{2}$-connected} \arrow[ur] \arrow[dr] &        \\
		&                                          & \textrm{strongly connected}
	\end{tikzcd}
\end{center}
	\caption[Relations among connectivity notions]{Implications for strong $2$-connectivity, strong $\frac{3}{2}$-connectivity, strong and weak connectivity. The converse implications do not hold, see \cref{prelim:fig:connectedness-notions}.}
	\label{prelim:fig:connectivity-notions}
\end{figure}

\begin{figure}[htb]
	\centering	
	\begin{subfigure}[b]{0.45\textwidth}
		\centering
		\begin{tikzpicture}
			\node[draw, circle] (1) {};
			\node[draw, circle, below left = of 1] (2) {};
			\node[draw, circle, below right = of 1] (3) {};
			
			\path[->] (1) edge (2)
			(2) edge (3)
			(3) edge (1);
		\end{tikzpicture}
		\caption{Strongly $\frac{3}{2}$-connected, but not strongly $2$-connected.}
	\end{subfigure}
	\hfill
	\begin{subfigure}[b]{0.45\textwidth}
		\centering
		\begin{tikzpicture}
			\node[draw, circle] (1) {};
			\node[draw, circle, right = of 1] (2) {};
			\node[draw, circle, right = of 2] (3) {};
			
			\path[->] (1) edge [bend right] (2)
			(2) edge [bend right] (3)
			(3) edge [bend right] (2)
			(2) edge [bend right] (1);
		\end{tikzpicture}
		\caption{Strongly connected, but not weakly $2$-connected.}
	\end{subfigure}
	
	\vspace{0.5cm}
	
	\begin{subfigure}[b]{0.45\textwidth}
		\centering
		\begin{tikzpicture}
			\node[draw, circle] (1) {};
			\node[draw, circle, below left = of 1] (2) {};
			\node[draw, circle, below right = of 1] (3) {};
			
			\path[->] (1) edge (2)
			(3) edge (2)
			(3) edge (1);
		\end{tikzpicture}
		\caption{Weakly $2$-connected, but not strongly connected.}
	\end{subfigure}
	\hfill
	\begin{subfigure}[b]{0.45\textwidth}
		\centering
		\begin{tikzpicture}
			\node[draw, circle] (1) {};
			\node[draw, circle, below = of 1] (2) {};
			\node[draw, circle, left = of 2] (3) {};
			\node[draw, circle, right = of 2] (4) {};
			
			\path[->] (1) edge (2)
			(2) edge (3)
			(2) edge (4)
			(3) edge (1)
			(4) edge (1);
		\end{tikzpicture}
		\caption{Strongly connected and weakly $2$-connected, but not strongly $\frac{3}{2}$-connected.}
	\end{subfigure}
	
	\caption[Examples for connectedness notions]{Some relations among the introduced notions of connectedness.}
	\label{prelim:fig:connectedness-notions}
\end{figure}

Lastly, we state a characterization of strong connectivity, which was given as a definition in~\cite[Def.~2]{DBLP:journals/jalc/GonzeJ19}.

\begin{lemma}
	\label{prelim:lem:partition_SCC}
	A digraph $G = (V, E)$ is strongly connected if and only if for
	every proper and non-empty $B \subsetneq V$
	there exists an edge $(u,v) \in E$ such that $u \notin B$ and $v \in B$. 
\end{lemma}
\begin{proof}
	First, suppose $G$ is strongly connected. Let $B \subsetneq V$
	with $u \in V \setminus B$ and~$v \in B$.
	By assumption there exists a path $v_0, v_1, \ldots, v_n$
	with $v_0 = u$ and $v_n = v$ in $G$.
	Let $i \in \{1,\ldots,n\}$
	be a largest index such that $v_{i} \notin B$.
	We have $i < n$ since $v_n \in B$, 
	and $v_{i+1} \in B$ by the maximal choice of~$i$.
	So, the edge $(v_i, v_{i+1}) \in E$
	fulfills the claim.
		
	Conversely, suppose the condition is true.
	Let $u,v \in V$ be distinct vertices.
	Set $P = \{ w \in V \mid \text{there exists a path from $u$ to $w$} \} \cup \{u\}$.
	If $V \ne P$, then, by applying the assumption to $B = V \setminus P$,
	there exists $(u', v') \in E$
	with $u' \in P$ and $v' \notin P$.
	As there exists a path from $u$ to $u'$,
	by adjoining $(u', v')$ we obtain a path from $u$ to $v'$.
	However, as $v' \notin P$ this is a contradiction.
	So, we must have $P = V$
	and there exists a path from $u$ to $v$.
	Hence, $G$ is strongly connected.
\end{proof} 

\begin{corollary}
	An undirected graph $G = (V,E)$ is connected
	if and only if for every proper and non-empty $B \subseteq V$
	there exists $e \in E$ with $|e \cap (V \setminus B)| = |e \cap B| = 1$.	
\end{corollary}
\begin{proof}
	Consider the digraph $(V, \{ (u,v), (v,u) \mid \{u,v\} \in E \})$.
	Since~$G$ is connected, it is strongly connected.
	By \cref{prelim:lem:partition_SCC} 
	it has an edge $(u,v)$ such that $u \notin B$ and $v \in B$.
	By definition $\{u,v\} \in E$.
	Then the claim holds for the edge $e = \{u,v\}$.
\end{proof}

\subsection{Graphs for Transformations and Homogeneous Relations}
\label{prelim:sec:graphs:funcgraph}

\label{def:GR} 
\label{def:Gf}
\IN{functional graph}
The \emph{functional graph} of a transformation $t\colon X \to X$
is the digraph~$G_t = (X, E)$ with 
\[
E = \{ (x, \applyfunc{x}{t}) \mid x \in X\}.
\]
In~$G_t$, the vertices in $\excl{t}$ are precisely those with indegree zero, and those in $\dupl{t}$ are precisely 
the vertices with indegree at least two. 
A digraph is the functional graph of some transformation
if and only if every vertex has outdegree one. 
These graphs are also called
maximal directed pseudoforests~\cite{GABOW1988259},
or sun-structures in~\cite{FernauGHHVW19}.

Furthermore, to every homogeneous relation~$R$ over~$X$
we associate a digraph 
$G_R = (X, E)$ where $E = \{ ( a, b ) \mid aRb \}$.

\begin{remark} \label{prelim:rem:intuition-functional-graph}
	The cycle points of a transformation are precisely those contained
	in cycles of the functional graph.
	The cycle points can be partitioned into points on individual cycles. The tail points
	are connected to the cycle points in a ``tree-like'' fashion 
	whose roots are precisely the root points.
	The edges on the tails point towards the root.
	See \cref{prelim:trans:fig:nonpermpoint} 
	where a functional graphs
	with two cycles and three tails is depicted.
	The roots are indicated by solid black circles.
		
	\begin{figure}[ht]
		\centering
		\begin{tikzpicture}
			\tikzset{every node/.style={draw,circle,inner sep=2.5pt}}
			
			\def\radius{4em}
			\def\shiftX{11em}
			\def\nA{10}
			\def\nB{8}
			
			\foreach \i in {1,2,4,5,7,8,9,10} {
				\node (A\i) at ({\radius*cos(360/\nA * \i)}, {\radius*sin(360/\nA * \i)}) {};
			}
			
			\foreach \i in {3,6} {
				\node [fill] (A\i) at ({\radius*cos(360/\nA * \i)}, {\radius*sin(360/\nA * \i)}) {};
			}
			
			\foreach \i in {2,...,\nA} {
				\path[->] (A\the\numexpr\i-1\relax) edge (A\i);
			} 
			\path[->] (A\nA) edge (A1);
			
			\foreach \j in {1,2,3,4,5,7,8} {
				\node (B\j) at ({\shiftX + 0.9*\radius*cos(360/\nB * \j)}, {1cm + 0.9*\radius*sin(360/\nB * \j)}) {};
			}
			
			\node [fill] (B6) at ({\shiftX + 0.9*\radius*cos(360/\nB *6)}, {1cm + 0.9*\radius*sin(360/\nB * 6)}) {};
			
			\foreach \i in {2,...,\nB} {
				\path[->] (B\the\numexpr\i-1\relax) edge (B\i);
			} 
			\path[->] (B\nB) edge (B1);
			
			\node[above = of A3] (T11) {};
			\node[left = of T11] (T12) {};
			
			\node[left = of A6] (T21) {};
			
			\node[below = of B6] (T31) {};
			\node[below left = of T31] (T32) {};
			\node[below right = of T31] (T33) {};
			
			\path[->] (T12) edge (T11); \path[->] (T11) edge (A3);
			\path[->] (T21) edge (A6);
			\path[->] (T33) edge (T31); 
			\path[->] (T32) edge (T31);
			\path[->] (T31) edge (B6);    
		\end{tikzpicture}
		\caption[Example functional graph]{Example of a functionl graph with root points marked, see \cref{prelim:rem:intuition-functional-graph}.}
		\label{prelim:trans:fig:nonpermpoint}	
	\end{figure}	
\end{remark}

\begin{proposition}
	\label{prelim:fact:relation_graph} Suppose $R$ is a homogeneous relation over $X$.
	Then the weakly connected components of $G_R$ 
	are precisely the equivalence classes of the equivalence relation generated by $R$.
\end{proposition}
\begin{proof}
	For $a, b \in X$
	we set $a \sim b$ 
	if $a$ and $b$ are contained in the same weakly connected component of $G_R$.
	Then $\sim$ is an equivalence relation such that, for $a,b \in X$, we have
	\begin{equation}\label{eqn:relation_graph}
	aRb \Rightarrow a \sim b.
	\end{equation}
		
	Now, let $E$ be any equivalence relation over $X$
	such that $a R b$ implies $a E b$ for all $a,b \in X$.
	Suppose $a \sim b$ for $a,b \in X$.
	Then there exist $v_0, v_1,\ldots, v_n \in X$ forming a path from $a$ to $b$ in the undirected graph
	resulting from $G_R$ by forgetting about the direction of edges, i.e., $v_0 = a$, $v_n = b$ and
	for each $i \in \{1,\ldots,n\}$ we have
	$
    v_{i-1} R v_i \text{ or } v_i R v_{i-1}.
    $
	Hence, by assumption and symmetry of $E$,
	we find $v_{i-1} E v_i$ for all $i \in \{1,\ldots,n\}$.
    By transitivity of $E$, we get $aEb$.
    Thus, $\sim$ refines~$E$.
    
    So, by \cref{eqn:relation_graph}
    and as $\sim$ refines every equivalence
    relation $E$ where $aRb$ implies $aEb$,
    it equals the equivalence relation generated by~$R$.
\end{proof}

For a more algebraic description of the equivalence relation
generated by~$R$, see~\cite[Prop.~1.4.9 \& 1.4.10]{howie1995fundamentals}.

\begin{remark}
	In the literature,
	the set $R_G \subseteq X \times Y$
	for a relation $R = (R_G, X, Y)$
	as defined in~\cref{def:relation}
	is often called the graph of the relation. However,
	since this is not a graph as defined
	here, we avoid this terminology.
\end{remark}